\documentclass[12pt,a4paper]{article}

% ── Packages ──
\usepackage[utf8]{inputenc}
\usepackage[T1]{fontenc}
\usepackage{lmodern}
\usepackage[british]{babel}
\usepackage[margin=2.5cm]{geometry}
\usepackage{graphicx}
\usepackage{booktabs}
\usepackage{array}
\usepackage{longtable}
\usepackage{tabularx}
\usepackage{hyperref}
\usepackage{listings}
\usepackage{xcolor}
\usepackage{float}
\usepackage{enumitem}
\usepackage{amsmath}
\usepackage{caption}
\usepackage{parskip}
\usepackage{fancyhdr}
\usepackage{titlesec}

% ── Colours ──
\definecolor{codebackground}{RGB}{245,245,245}
\definecolor{codekeyword}{RGB}{0,0,180}
\definecolor{codecomment}{RGB}{0,128,0}
\definecolor{codestring}{RGB}{163,21,21}
\definecolor{linkblue}{RGB}{0,70,140}

% ── Hyperref ──
\hypersetup{
    colorlinks=true,
    linkcolor=linkblue,
    citecolor=linkblue,
    urlcolor=linkblue
}

% ── Listings ──
\lstset{
    backgroundcolor=\color{codebackground},
    basicstyle=\ttfamily\small,
    keywordstyle=\color{codekeyword}\bfseries,
    commentstyle=\color{codecomment},
    stringstyle=\color{codestring},
    breaklines=true,
    frame=single,
    rulecolor=\color{black!30},
    numbers=none,
    xleftmargin=4pt,
    xrightmargin=4pt,
    aboveskip=10pt,
    belowskip=10pt,
    showstringspaces=false,
    tabsize=2
}

\lstdefinelanguage{json}{
    morestring=[b]",
    morecomment=[l]{//},
    literate=
      {0}{{{\color{codekeyword}0}}}{1}
      {1}{{{\color{codekeyword}1}}}{1}
      {2}{{{\color{codekeyword}2}}}{1}
      {3}{{{\color{codekeyword}3}}}{1}
      {4}{{{\color{codekeyword}4}}}{1}
      {5}{{{\color{codekeyword}5}}}{1}
      {6}{{{\color{codekeyword}6}}}{1}
      {7}{{{\color{codekeyword}7}}}{1}
      {8}{{{\color{codekeyword}8}}}{1}
      {9}{{{\color{codekeyword}9}}}{1}
      {:}{{{\color{black}{:}}}}{1}
      {,}{{{\color{black}{,}}}}{1}
}

% ── Headers ──
\pagestyle{fancy}
\fancyhf{}
\fancyhead[L]{\small Loki IDS}
\fancyhead[R]{\small Security Engineering Lab}
\fancyfoot[C]{\thepage}
\renewcommand{\headrulewidth}{0.4pt}

% ── Title spacing ──
\titleformat{\section}{\Large\bfseries}{Task \thesection:\ }{0em}{}
\titleformat{\subsection}{\large\bfseries}{\thesubsection\ }{0em}{}
\titleformat{\subsubsection}{\normalsize\bfseries}{\thesubsubsection\ }{0em}{}

% ══════════════════════════════════════════════════════════
\begin{document}

% ── Title Page ──
\begin{titlepage}
    \centering
    \vspace*{2cm}

    {\large Lehrstuhl f\"ur Technische Informatik\\
    Prof.\ Dr.\ Stefan Katzenbeisser\\[4pt]
    Organiser: Martin Schmid\\
    Supervisor: Emiliia Geloczi}

    \vspace{1.5cm}
    {\large Security Engineering Lab, Advanced Security Engineering Lab}

    \vspace{1cm}
    \rule{\textwidth}{1pt}
    \vspace{0.5cm}

    {\Huge \textbf{Loki IDS:\\Intrusion Detection in Smart Home}}

    \vspace{0.5cm}
    \rule{\textwidth}{1pt}

    \vspace{2cm}
    {\large Final Report}

    \vspace{1.5cm}
    {\large
    Omar Shehata \\
    Ahmed Ghanem \\
    Moustafa Elgendy \\
    Mohammed Zaher}

    \vfill
    {\large Universit\"at Passau \\ Winter Semester 2025/26}
\end{titlepage}

% ── Table of Contents ──
\newpage
\tableofcontents
\newpage

% ══════════════════════════════════════════════════════════
% TASK 1 — THEORETICAL BACKGROUND
% ══════════════════════════════════════════════════════════
\section{Theoretical Background about IoT Systems}
\label{sec:theory}

\subsection{Internet of Things (IoT)}
\label{sec:iot}

The Internet of Things (IoT) refers to a network of physical objects---sensors, appliances, vehicles, and other everyday devices---embedded with electronics, software, and connectivity that allows them to collect and exchange data over the internet or local networks. Unlike traditional computing devices, IoT endpoints are typically resource-constrained (limited CPU, memory, and battery) and designed for a specific purpose, such as measuring temperature, detecting motion, or controlling a light bulb.

In a smart home context, IoT devices include smart locks, thermostats, cameras, motion sensors, and lighting systems. These devices communicate with a central hub or gateway, which coordinates their behaviour and provides user-facing control interfaces. The value lies in automation and remote control, but the always-on, always-connected nature of these devices also means they continuously generate and transmit data, making them attractive targets for adversaries~\cite{atzori2010iot, alfuqaha2015iot}.

\subsection{Architecture of IoT Systems}
\label{sec:architecture}

IoT systems are commonly described using a layered architecture. Whilst the exact number of layers varies across the literature, the most widely adopted model consists of three to five layers~\cite{khan2012future, wu2010architecture}:

\subsubsection*{Three-Layer Model}

\begin{description}[style=nextline]
    \item[Perception Layer (Device Layer)]
    The physical layer where sensors and actuators interact with the environment. Devices here collect raw data (e.g.\ temperature readings, motion events) and execute commands (e.g.\ turn on an LED, sound a buzzer). Common hardware includes microcontrollers such as ESP32 or Arduino boards. Communication at this layer typically uses low-power protocols like Bluetooth Low Energy (BLE), Zigbee, or direct Wi-Fi.

    \item[Network Layer (Transport Layer)]
    Responsible for transmitting data between the perception layer and the application layer. This layer handles routing, addressing, and protocol translation. In smart homes, Wi-Fi is the dominant transport medium, and protocols such as MQTT (Message Queuing Telemetry Transport), CoAP (Constrained Application Protocol), and HTTP/HTTPS are used for message exchange. MQTT is particularly common in IoT due to its lightweight publish-subscribe model, which is well suited for constrained devices~\cite{banks2014mqtt}.

    \item[Application Layer]
    Provides the user-facing services and interfaces. This includes dashboards, mobile applications, and automation engines that consume data from the lower layers and present it in a meaningful form. In our system, this corresponds to the web-based dashboard served by the Raspberry~Pi.
\end{description}

\begin{table}[H]
\centering
\caption{Key protocols used in our IoT system.}
\label{tab:protocols}
\begin{tabular}{lll}
\toprule
\textbf{Layer} & \textbf{Protocol} & \textbf{Purpose} \\
\midrule
Perception  & GPIO, I2C           & Sensor/actuator interfacing \\
Network     & Wi-Fi (802.11)      & Wireless connectivity \\
Network     & MQTT v3.1.1         & Lightweight publish-subscribe messaging \\
Network     & HTTP/REST           & API communication \\
Application & WebSocket           & Real-time dashboard updates \\
\bottomrule
\end{tabular}
\end{table}

\subsection{Vulnerabilities of IoT Systems}
\label{sec:vulnerabilities}

IoT devices suffer from a number of inherent vulnerabilities that stem from their constrained nature and the way they are typically deployed~\cite{owasp2018iot, bertino2017botnets}:

\begin{itemize}
    \item \textbf{Weak or default credentials:} Many IoT devices ship with factory-set passwords that users never change, making brute-force or dictionary attacks trivial.
    \item \textbf{Lack of encryption:} Resource constraints often lead manufacturers to omit TLS or use outdated cryptographic primitives, exposing data in transit to eavesdropping.
    \item \textbf{Insufficient update mechanisms:} Many IoT devices lack support for over-the-air (OTA) firmware updates, meaning known vulnerabilities remain unpatched indefinitely.
    \item \textbf{Large attack surface:} Each connected device is a potential entry point. A single compromised sensor can serve as a pivot to reach other devices on the same network.
    \item \textbf{Insecure network services:} Devices may expose unnecessary ports or run services with known flaws (e.g.\ Telnet, UPnP).
    \item \textbf{Limited computational resources:} Constrained memory and processing power make it difficult to run traditional security software directly on the devices.
    \item \textbf{Physical accessibility:} IoT devices are often deployed in accessible locations, making them susceptible to physical tampering.
\end{itemize}

\subsection{Existing Attacks Against IoT Systems}
\label{sec:attacks}

The following are well-documented attack categories relevant to IoT environments~\cite{ferrag2020deep, neshenko2019demystifying, antonakakis2017mirai}:

\subsubsection*{Network-Layer Attacks}

\begin{description}[style=nextline]
    \item[Port Scanning]
    An attacker probes a range of ports on a target device to discover open services. Tools such as Nmap automate this process. In IoT networks, port scans are often a precursor to more targeted exploitation.

    \item[Denial of Service (DoS/DDoS)]
    The attacker floods a target with traffic to exhaust its resources and render it unresponsive. Common variants include:
    \begin{itemize}
        \item \emph{TCP SYN Flood:} Sends a high volume of TCP SYN packets without completing the handshake, consuming the target's connection table.
        \item \emph{UDP Flood:} Sends large volumes of UDP datagrams to random ports, forcing the target to process and respond to each one.
        \item \emph{ICMP Flood (Ping Flood):} Overwhelms a target with ICMP echo request packets.
    \end{itemize}

    \item[Man-in-the-Middle (MitM)]
    The attacker intercepts communication between two parties (e.g.\ between an ESP32 device and the MQTT broker) by positioning themselves in the network path, allowing eavesdropping and message manipulation.

    \item[ARP Spoofing]
    The attacker sends forged ARP messages to associate their MAC address with the IP address of a legitimate device, redirecting traffic through their machine.
\end{description}

\subsubsection*{Application-Layer Attacks}

\begin{description}[style=nextline]
    \item[Replay Attacks]
    An attacker captures legitimate MQTT messages and re-transmits them to trigger unintended actions (e.g.\ replaying an ``unlock door'' command).

    \item[Packet Injection]
    Crafting and injecting malicious packets into the network to exploit vulnerabilities in application-layer parsers. This includes patterns like SQL injection attempts or cross-site scripting (XSS) payloads embedded in HTTP traffic.
\end{description}

\subsubsection*{Device-Level Attacks}

\begin{description}[style=nextline]
    \item[Firmware Extraction and Reverse Engineering]
    Physical access to the device allows extraction of firmware via JTAG or UART interfaces, potentially revealing hardcoded credentials or cryptographic keys.

    \item[Botnet Recruitment]
    Compromised IoT devices are enrolled into botnets (e.g.\ Mirai) and used to launch large-scale DDoS attacks against third-party targets~\cite{antonakakis2017mirai}.
\end{description}

\subsection{Literature Review (PRISMA)}
\label{sec:prisma}

A literature review following the PRISMA (Preferred Reporting Items for Systematic Reviews and Meta-Analyses) methodology~\cite{page2021prisma} was conducted to survey existing intrusion detection systems designed for IoT environments.

\subsubsection*{Search Strategy}

Databases searched included IEEE Xplore, ACM Digital Library, and ScienceDirect. The search terms used were: (\emph{``intrusion detection system''} OR \emph{``IDS''}) AND (\emph{``IoT''} OR \emph{``Internet of Things''} OR \emph{``smart home''}). The search was limited to publications from 2018 to 2024.

\subsubsection*{Inclusion Criteria}
\begin{itemize}
    \item Peer-reviewed journal articles or conference papers.
    \item Proposed or implemented an IDS specifically for IoT or smart home environments.
    \item Written in English.
\end{itemize}

\subsubsection*{Exclusion Criteria}
\begin{itemize}
    \item Survey-only papers without a concrete IDS proposal.
    \item Papers focusing exclusively on industrial IoT (IIoT) or vehicular networks.
    \item Duplicate publications.
\end{itemize}

\begin{table}[H]
\centering
\caption{PRISMA screening flow.}
\label{tab:prisma}
\begin{tabular}{lr}
\toprule
\textbf{Stage} & \textbf{Count} \\
\midrule
Records identified through database search & $\sim$320 \\
After removing duplicates                  & $\sim$210 \\
After title/abstract screening             & $\sim$65 \\
After full-text assessment                 & 18 \\
Studies included in final review           & 18 \\
\bottomrule
\end{tabular}
\end{table}

The 18 selected studies were analysed for their detection approach, deployment location, detected attack types, and reported performance metrics.

\subsection{Existing IDSs: Limitations and Capabilities}
\label{sec:ids_review}

Based on the literature review, the identified IDS solutions can be grouped into three main categories:

\subsubsection*{1. Signature-Based IDSs}
\begin{itemize}
    \item \emph{Capabilities:} High detection accuracy for known attack patterns. Low false-positive rate when signatures are well defined. Straightforward to implement and understand.
    \item \emph{Limitations:} Cannot detect novel or zero-day attacks. Require continuous signature updates. Large signature databases consume memory, which is problematic on constrained devices.
    \item \emph{Examples:} Snort-based approaches adapted for IoT gateways~\cite{khraisat2019survey}, lightweight rule-matching engines on Raspberry~Pi platforms.
\end{itemize}

\subsubsection*{2. Anomaly-Based IDSs (Machine Learning)}
\begin{itemize}
    \item \emph{Capabilities:} Can detect previously unseen attacks by modelling normal behaviour. Adaptable to different network environments without manual rule creation.
    \item \emph{Limitations:} High false-positive rates, especially during the training phase. Require substantial computational resources for model training and inference. Difficult to interpret and debug (black-box models). Sensitive to concept drift when normal traffic patterns change over time~\cite{sommer2010outside}.
    \item \emph{Examples:} Random Forest classifiers on network flow features, autoencoders for anomaly detection, federated learning approaches for distributed IoT networks.
\end{itemize}

\subsubsection*{3. Hybrid IDSs}
\begin{itemize}
    \item \emph{Capabilities:} Combine signature and anomaly detection to achieve broader coverage. Can detect known attacks precisely whilst still flagging anomalous behaviour.
    \item \emph{Limitations:} Increased complexity and resource consumption. Integration of two detection paradigms introduces design challenges. May still suffer from the individual weaknesses of each approach.
    \item \emph{Examples:} Systems that use signature matching as a first pass and fall back to statistical analysis for unmatched traffic~\cite{alaiz2019multiclass}.
\end{itemize}

\subsubsection*{Research Gap Identified}

Most existing solutions either require significant computational resources (making them unsuitable for a Raspberry~Pi gateway) or focus solely on one detection paradigm. Few solutions offer real-time monitoring with a practical web-based dashboard and integrated IoT device control. Additionally, many lack proper alert lifecycle management---tracking when an attack starts, progresses, and ends. Loki~IDS addresses these gaps by combining lightweight behavioural detection with signature matching, deployed directly on the gateway, with full alert lifecycle tracking and an integrated device control interface.


% ══════════════════════════════════════════════════════════
% TASK 2 — TEST ENVIRONMENT ESTABLISHMENT
% ══════════════════════════════════════════════════════════
\newpage
\section{Test Environment Establishment}
\label{sec:environment}

This section describes the physical and software environment built to develop and evaluate the IDS. The test-bed consists of a Raspberry~Pi~5 acting as a Wi-Fi access point, network gateway, MQTT broker, and API server, with two ESP32 microcontrollers representing typical smart home IoT devices. A second Raspberry~Pi serves as the attacker machine, connected to the network through the MikroTik router rather than through the access point, simulating an external adversary.

\subsection{Hardware Components}
\label{sec:hardware}

Table~\ref{tab:hardware} lists all hardware used in the test environment.

\begin{table}[H]
\centering
\caption{Hardware inventory.}
\label{tab:hardware}
\begin{tabular}{llp{6.5cm}}
\toprule
\textbf{Component} & \textbf{Model} & \textbf{Role} \\
\midrule
Central gateway  & Raspberry Pi 5 (8\,GB) & Access point, router, MQTT broker, IDS host, API server \\
IoT device 1     & ESP32 DevKit v1        & Entryway security hub (PIR sensor + buzzer + LED) \\
IoT device 2     & ESP32 DevKit v1        & Smart bulb controller (NeoPixel RGB LED strip) \\
Attacker machine & Raspberry Pi           & Attack simulation and IDS testing \\
Network uplink   & MikroTik Router        & Internet connectivity and attacker network via Ethernet \\
\bottomrule
\end{tabular}
\end{table}

\subsection{Network Architecture}
\label{sec:network}

The Raspberry~Pi creates an isolated Wi-Fi network to which all IoT devices and management clients connect. The Pi's wireless interface (\texttt{wlan0}) is assigned the static IP \texttt{10.0.0.1/24} and acts as the DHCP server for the subnet \texttt{10.0.0.0/24}. The Ethernet interface (\texttt{eth0}) connects to a MikroTik router, which provides internet access and serves as the network entry point for the attacker machine. The attacker Raspberry~Pi connects through the router---not through the access point---simulating an external adversary who has network-level reachability to the gateway but is not a member of the IoT Wi-Fi network. Attacks arriving via the router enter the Pi through \texttt{eth0} and are intercepted on the INPUT chain (traffic destined to the Pi) or the FORWARD chain (traffic routed towards IoT devices on \texttt{wlan0}).

\begin{lstlisting}[frame=none,backgroundcolor=\color{white}]
                      Internet
                         |
                  [MikroTik Router]
                   /             \
           eth0  /               \
   [Raspberry Pi AP]        [Attacker RPi]
    wlan0: 10.0.0.1/24
       /       |       \
 [ESP32-1]  [ESP32-2]  [Browser]
 10.0.0.x   10.0.0.x   10.0.0.x
\end{lstlisting}

\subsubsection*{Raspberry Pi Roles}

The Pi serves four simultaneous roles:

\begin{enumerate}
    \item \textbf{Wireless Access Point (AP):} Configured using RaspAP (which wraps \texttt{hostapd} and \texttt{dnsmasq}) to broadcast a Wi-Fi network named \texttt{Access-Point}. Legacy \texttt{iptables} is selected via \texttt{update-alternatives} for compatibility with NFQUEUE.

    \item \textbf{Network Gateway / Router:} IP forwarding is enabled (\texttt{sysctl -w net.ipv4.ip\_forward=1}) so that traffic from IoT devices can be routed through the Pi. Selective port forwarding rules route HTTP/HTTPS traffic from \texttt{eth0} to an internal host when needed:
    \begin{lstlisting}[language=bash,frame=none,backgroundcolor=\color{white}]
sudo iptables -t nat -A PREROUTING -i eth0 -p tcp \
     --dport 80 -j DNAT --to-destination 10.0.0.6
sudo iptables -I FORWARD -i eth0 -o wlan0 -p tcp \
     --dport 80 -d 10.0.0.6 -j ACCEPT
    \end{lstlisting}

    \item \textbf{MQTT Broker:} Mosquitto listens on port 1883 and acts as the central message broker for all publish-subscribe communication between ESP32 devices and the dashboard.

    \item \textbf{IoT Hub / API Server:} A FastAPI-based web server runs on port 8080, serving the REST API, the web dashboard, and relaying MQTT commands to devices.
\end{enumerate}

\subsubsection*{Static IP Configuration}

To ensure the Pi's network interfaces remain statically configured, the Ethernet interface (\texttt{eth0}) is excluded from DHCP management in \texttt{dhcpcd.conf}:

\begin{lstlisting}[language=bash,frame=none,backgroundcolor=\color{white}]
# /etc/dhcpcd.conf
denyinterfaces eth0

# systemd-networkd (alternative)
[Match]
Name=wlan0

[Network]
Address=10.0.0.1/24
DHCP=no
IPForward=yes
\end{lstlisting}

\subsection{IoT Device Selection and Functionality}
\label{sec:devices}

Based on the available hardware (two ESP32 modules with sensor kits), we designed two IoT devices that represent common smart home components:

\subsubsection*{ESP32-1: Entryway Security Hub (Motion Sensor + Alarm)}
\begin{itemize}
    \item \textbf{Sensors/Actuators:} PIR motion sensor (GPIO~15), piezo buzzer (GPIO~3), status LED (GPIO~2).
    \item \textbf{Functionality:} Continuously monitors for motion via a polling loop. When motion is detected, the device publishes a JSON alert over MQTT to \texttt{esp32/sensor1/status} and triggers the buzzer with a two-tone alarm pattern. The alarm automatically resets after 5~seconds of inactivity. Additionally, upon detecting motion, ESP32-1 sends a command to ESP32-2 to turn on the RGB LEDs at maximum brightness, demonstrating inter-device coordination. All components (alarm, buzzer, LED) can be independently controlled from the web dashboard via MQTT commands.
    \item \textbf{Dashboard commands:} \texttt{alarm\_control} (enable/disable/test), \texttt{buzzer\_control} (on/off/beep with configurable duration), \texttt{led\_control} (on/off/auto).
    \item \textbf{Use case:} Simulates a smart home entryway sensor that detects intruders and raises an audible alarm.
\end{itemize}

\subsubsection*{ESP32-2: Smart Bulb Controller (RGB LED)}
\begin{itemize}
    \item \textbf{Sensors/Actuators:} NeoPixel RGB LED strip (4~LEDs on GPIO~15) driven by the \texttt{Adafruit\_NeoPixel} library.
    \item \textbf{Functionality:} Controls an RGB LED strip with on/off switching and brightness adjustment (0--255). Receives \texttt{bulb\_control} commands from the dashboard (or from ESP32-1) via MQTT on the \texttt{rpi/broadcast} topic. Publishes its current state---including heartbeat messages every 30~seconds---to \texttt{esp32/sensor2/status}.
    \item \textbf{Use case:} Simulates a smart light bulb that can be remotely controlled from the dashboard.
\end{itemize}

\subsection{Communication Protocol: MQTT}
\label{sec:mqtt}

All device-to-gateway communication uses the MQTT v3.1.1 publish-subscribe protocol~\cite{banks2014mqtt}. The Mosquitto broker on the Pi mediates all messages. Devices subscribe to \texttt{rpi/broadcast} for commands and publish status on dedicated topics.

\begin{table}[H]
\centering
\caption{MQTT topic structure.}
\label{tab:mqtt_topics}
\begin{tabular}{llp{5.5cm}}
\toprule
\textbf{Topic} & \textbf{Direction} & \textbf{Purpose} \\
\midrule
\texttt{rpi/broadcast}         & Pi $\rightarrow$ Devices & Dashboard commands (JSON) \\
\texttt{esp32/sensor1/status}  & ESP32-1 $\rightarrow$ Pi & Motion events, alarm state, heartbeats \\
\texttt{esp32/sensor2/status}  & ESP32-2 $\rightarrow$ Pi & Bulb state, heartbeats \\
\bottomrule
\end{tabular}
\end{table}

\subsubsection*{Message Format Examples}

Command from dashboard to ESP32-1:
\begin{lstlisting}[language=json]
{
  "device": "esp32-1",
  "command": "alarm_control",
  "action": "enable"
}
\end{lstlisting}

Status message from ESP32-1 to dashboard:
\begin{lstlisting}[language=json]
{
  "device": "esp32-1",
  "event": "motion_detected",
  "alarm_enabled": true,
  "alarm_active": true,
  "buzzer_on": false,
  "led_auto_mode": true
}
\end{lstlisting}

Command from dashboard to ESP32-2:
\begin{lstlisting}[language=json]
{
  "device": "esp32-2",
  "command": "bulb_control",
  "state": "on",
  "brightness": 200
}
\end{lstlisting}

\subsection{ESP32 Firmware}
\label{sec:firmware}

Both ESP32 devices run firmware written in C using the Arduino framework. The firmware handles Wi-Fi connectivity, MQTT communication (via the \texttt{PubSubClient} library), JSON parsing (via \texttt{ArduinoJson}), and sensor/actuator control. Both devices implement automatic Wi-Fi and MQTT reconnection if the connection drops, and use non-blocking reconnect logic with 5-second retry intervals.

\subsubsection*{ESP32-1 Firmware (\texttt{esp/full-code.ino})}
\begin{itemize}
    \item Connects to the Pi's Wi-Fi AP (\texttt{Access-Point}) and establishes an MQTT session with the broker at \texttt{10.0.0.1:1883}.
    \item Subscribes to \texttt{rpi/broadcast} for receiving commands from the dashboard.
    \item Reads the PIR sensor on GPIO~15 using a polling loop in \texttt{loop()}.
    \item On motion detection, publishes a JSON status message and triggers a two-tone buzzer alarm. Simultaneously sends a \texttt{bulb\_control} command to ESP32-2 to turn on the RGB LEDs.
    \item Supports three command types: \texttt{alarm\_control}, \texttt{buzzer\_control}, and \texttt{led\_control}.
\end{itemize}

\subsubsection*{ESP32-2 Firmware (\texttt{esp/Final-RGB.ino})}
\begin{itemize}
    \item Connects to the same AP and MQTT broker.
    \item Controls a NeoPixel RGB LED strip (4~LEDs) on GPIO~15 using the \texttt{Adafruit\_NeoPixel} library.
    \item Responds to \texttt{bulb\_control} commands with on/off states and brightness levels from 0 to 255.
    \item Sends periodic heartbeat messages (every 30~seconds) to indicate it is online.
\end{itemize}

\subsection{Web Interface for Device Control}
\label{sec:webinterface}

The web interface is a single-page dashboard built with HTML, CSS, and vanilla JavaScript (no frameworks). It is served by the FastAPI backend on port 8080 and provides four main views:

\begin{itemize}
    \item \textbf{Dashboard Tab:} Displays real-time statistics---total alerts, alerts in the last 24~hours, IDS running status, a bar chart of alerts by type (Signature, Behaviour, System), and a list of the top attacking IP addresses.
    \item \textbf{Alerts Tab:} Lists all alerts with filtering by type, subtype, status, pattern, source IP, and destination IP. Supports pagination and shows the full alert lifecycle (STARTED, ONGOING, ENDED).
    \item \textbf{Signatures Tab:} Allows creation, editing, deletion, and import (from YAML files) of detection signatures. Each signature has a name, pattern, description, and enabled/disabled status.
    \item \textbf{IoT Control Tab:} Shows all registered ESP32 devices with their online/offline status (determined by MQTT heartbeats). Provides device-specific controls: alarm enable/disable/test, buzzer on/off/beep, LED on/off/auto for ESP32-1, and bulb on/off with brightness slider for ESP32-2.
\end{itemize}

The dashboard uses Chart.js for visualisations and connects to a WebSocket endpoint (\texttt{/ws/alerts}) for real-time alert notifications via toast messages.

\subsection{Software Stack}
\label{sec:software_stack}

Table~\ref{tab:software} summarises the software components used in the test environment.

\begin{table}[H]
\centering
\caption{Software stack.}
\label{tab:software}
\begin{tabular}{llp{5cm}}
\toprule
\textbf{Component} & \textbf{Technology} & \textbf{Purpose} \\
\midrule
Operating System    & Raspberry Pi OS (Debian) & Base OS on Pi \\
Access Point        & RaspAP (hostapd + dnsmasq) & Wi-Fi AP and DHCP \\
MQTT Broker         & Mosquitto 2.x            & Pub-sub messaging \\
API Server          & FastAPI + Uvicorn         & REST API and static file serving \\
Database            & SQLite (WAL mode)         & Alert and signature storage \\
Packet Interception & iptables + NFQUEUE        & Kernel-level packet capture \\
Packet Parsing      & Scapy + NetfilterQueue    & Userspace packet analysis \\
MQTT Client (Pi)    & paho-mqtt 2.x (Python)    & Dashboard $\leftrightarrow$ device bridge \\
MQTT Client (ESP32) & PubSubClient (C++)        & Device-side MQTT \\
JSON Parsing (ESP32)& ArduinoJson              & Command parsing on devices \\
LED Control         & Adafruit NeoPixel         & RGB strip driver on ESP32-2 \\
Frontend            & HTML + CSS + JS + Chart.js & Web dashboard \\
Python Environment  & Python 3.13 (venv)        & Virtual environment for all dependencies \\
\bottomrule
\end{tabular}
\end{table}

\subsection{Reproducible Setup Steps}
\label{sec:env_steps}

\subsubsection*{Prerequisites}
\begin{itemize}
    \item Raspberry~Pi 5 with Raspberry~Pi OS (Desktop).
    \item Two ESP32 DevKit modules with sensor kits.
    \item Python 3.12 or later installed on the Pi (Python 3.13 recommended).
    \item Arduino IDE or PlatformIO for ESP32 firmware flashing.
\end{itemize}

\subsubsection*{Step 1: Configure the Raspberry Pi as an Access Point}

\begin{lstlisting}[language=bash]
# Install dependencies and RaspAP
sudo apt update && sudo apt upgrade -y
sudo apt install iptables -y
sudo update-alternatives --set iptables /usr/sbin/iptables-legacy
sudo update-alternatives --set ip6tables /usr/sbin/ip6tables-legacy
curl -sL https://install.raspap.com | bash
sudo reboot

# Verify Wi-Fi is unblocked
rfkill list
rfkill unblock wifi
sudo systemctl restart hostapd

# Prevent wlan0 from getting DHCP from upstream router
sudo nano /etc/dhcpcd.conf
# Add: denyinterfaces eth0

# Ensure wlan0 has static IP 10.0.0.1/24 after each reboot
sudo systemctl restart dhcpcd
\end{lstlisting}

\subsubsection*{Step 2: Install Mosquitto MQTT Broker}

\begin{lstlisting}[language=bash]
sudo apt install mosquitto mosquitto-clients -y
sudo systemctl enable mosquitto
sudo systemctl start mosquitto
\end{lstlisting}

\subsubsection*{Step 3: Clone the Repository and Install Dependencies}

\begin{lstlisting}[language=bash]
cd ~
git clone <repository_url> Loki-IDS
cd Loki-IDS
bash setup_venv.sh
# NetfilterQueue system dependency:
sudo apt install python3-dev libnetfilter-queue-dev
\end{lstlisting}

This creates a virtual environment at \texttt{loki\_env/} and installs all Python dependencies (FastAPI, Uvicorn, SQLAlchemy, Scapy, NetfilterQueue, paho-mqtt, etc.).

\subsubsection*{Step 4: Flash ESP32 Firmware}

\begin{enumerate}
    \item Open Arduino IDE (or PlatformIO).
    \item Install the ESP32 board package and the following libraries: \texttt{PubSubClient}, \texttt{Adafruit\_NeoPixel}, \texttt{ArduinoJson}.
    \item Open \texttt{esp/full-code.ino}, update the Wi-Fi SSID and password to match the AP configuration, and flash it to ESP32-1.
    \item Open \texttt{esp/Final-RGB.ino}, update the Wi-Fi credentials, and flash it to ESP32-2.
\end{enumerate}

\subsubsection*{Step 5: Start the System}

\begin{lstlisting}[language=bash]
sudo bash start_loki_system.sh
\end{lstlisting}

This script starts the API server on port 8080, waits for it to become healthy, sets up iptables NFQUEUE rules, and launches the IDS engine. The dashboard is then accessible at \texttt{http://10.0.0.1:8080}.

\subsubsection*{Step 6: Verify}

\begin{lstlisting}[language=bash]
curl http://localhost:8080/api/system/health
curl http://localhost:8080/api/system/status
mosquitto_sub -t "esp32/+/status" -v
\end{lstlisting}


% ══════════════════════════════════════════════════════════
% TASK 3 — IMPLEMENTATION OF IDS
% ══════════════════════════════════════════════════════════
\newpage
\section{Implementation of IDS}
\label{sec:ids}

\subsection{Introduction and Motivation}
\label{sec:ids_intro}

Loki IDS is a lightweight, network-based Intrusion Detection System designed to protect smart home IoT environments. It is deployed on a Raspberry~Pi that acts simultaneously as a Wi-Fi access point, IoT gateway, and MQTT broker. All network traffic flowing through the gateway is intercepted and inspected in real time, enabling the detection of both known attack signatures and anomalous behavioural patterns.

The system employs a hybrid detection strategy that combines signature-based pattern matching with statistical behavioural analysis. The behavioural engine uses a dual-threshold approach: a sliding window counts the number of packets within a fixed time interval, whilst an Exponentially Weighted Moving Average (EWMA) estimator provides a smoothed, real-time measurement of the packet arrival rate in packets per second.

Crucially, the EWMA thresholds are \emph{adaptive}. Rather than relying on fixed, manually-tuned rate limits, each network flow learns its own baseline ``normal'' rate using a slow-moving EWMA. The detection threshold is then computed as a configurable multiple of that baseline (e.g.\ $5\times$ normal). This means the IDS automatically adapts to the actual traffic profile of the network: a flow that normally sees 10~packets per second will have a threshold of 50~pps, whilst a flow that normally sees 2~pps will have a threshold of 10~pps. A minimum floor value ensures the system remains protected during the initial warm-up period before a reliable baseline has been established.

Loki IDS is capable of detecting the following attack categories:
\begin{itemize}
    \item TCP SYN port scanning,
    \item TCP SYN flood (DoS/DDoS),
    \item UDP flood (DoS/DDoS),
    \item ICMP flood (DoS/DDoS), and
    \item arbitrary payload-based attacks via configurable signature rules (e.g.\ SQL injection, XSS, path traversal).
\end{itemize}

All alerts are logged to structured JSONL files and concurrently written to an SQLite database for integration with the web dashboard.

\subsection{System and Adversary Model}
\label{sec:models}

\subsubsection*{System Model}

The system consists of a Raspberry~Pi acting as the central gateway for a small smart home network. All IoT traffic---both between devices and between devices and the internet---passes through the Pi. The Pi runs three main services:

\begin{enumerate}
    \item \textbf{Network services:} Access point (\texttt{hostapd}), DHCP (\texttt{dnsmasq}), routing (iptables with IP forwarding).
    \item \textbf{MQTT broker:} Mosquitto, handling all publish-subscribe communication between ESP32 devices and the dashboard.
    \item \textbf{Loki IDS and API server:} The intrusion detection engine and the web interface, both running on the Pi.
\end{enumerate}

The IDS is positioned inline on the gateway, intercepting packets on both the INPUT chain (traffic destined to the Pi itself) and the FORWARD chain (traffic passing through the Pi between IoT devices and the external network). This placement gives the IDS visibility over all network traffic without requiring any software on the IoT devices themselves.

\subsubsection*{Adversary Model}

We consider an adversary with the following characteristics:

\begin{itemize}
    \item \textbf{Network access:} The adversary is either connected to the same Wi-Fi network as the IoT devices (insider threat) or is on an adjacent network reachable through the router (external attacker).
    \item \textbf{Capabilities:} The adversary can send arbitrary packets to any device on the network. They can perform port scans, launch flood-based DoS attacks, and inject malicious payloads into network traffic.
    \item \textbf{Goals:} The adversary aims to disrupt IoT device operation (availability attacks), discover running services for exploitation (reconnaissance), or inject malicious content into network traffic (integrity attacks).
    \item \textbf{Limitations:} The adversary does not have physical access to the Raspberry~Pi or ESP32 devices. The adversary cannot compromise the IDS itself (the IDS is trusted). The adversary does not have credentials for the Mosquitto broker (if authentication is enabled).
\end{itemize}

\subsection{System Architecture Overview}
\label{sec:arch_overview}

The Loki IDS is composed of several modular components that operate together in a pipeline architecture. Figure~\ref{fig:arch} illustrates the data flow from packet arrival to alert generation.

\begin{figure}[H]
\begin{lstlisting}[frame=none,backgroundcolor=\color{white},basicstyle=\ttfamily\footnotesize]
  Network Traffic (Wi-Fi AP)
           |
           v
  [iptables NFQUEUE] --- INPUT chain  (queue 100)
           |              FORWARD chain (queue 200)
           v
  [Packet Parser] --- packet_parser.py
  (Scapy: IP/TCP/UDP/ICMP extraction)
           |
    +------+------+
    |             |
    v             v
 [Behaviour    [Signature
  Engine]       Engine]
  detectore_    signature_
  engine.py     engine.py
 (Sliding       (Payload pattern
  Window +       matching from DB)
  Adaptive
  EWMA)
    |             |
    +------+------+
           |
           v
  [Logger Module] --- logger.py
  (Deduplication, Lifecycle:
   STARTED -> ONGOING -> ENDED)
           |
    +------+------+
    |             |
    v             v
 [JSONL File]  [SQLite DB]
  logs/         (via db_integration.py)
\end{lstlisting}
\caption{Loki IDS data flow architecture.}
\label{fig:arch}
\end{figure}

The system runs two concurrent threads: one bound to NFQUEUE~100 (INPUT chain, traffic destined for the Raspberry~Pi itself) and one bound to NFQUEUE~200 (FORWARD chain, traffic passing through the Pi to IoT devices). Both threads share the same signature engine instance but maintain independent behavioural detection state to avoid cross-chain interference.

\subsection{Network Interception Layer}
\label{sec:interception}

\subsubsection{iptables and NFQUEUE}

Loki IDS uses the Linux Netfilter framework to intercept packets in the kernel before they are delivered to their destination. The \texttt{iptables} utility inserts NFQUEUE rules that redirect packets to userspace queues where the Python application can inspect them. The \texttt{-{}-queue-bypass} flag ensures that network traffic continues to flow even if the IDS process is not running or crashes.

The following iptables rules are applied by \texttt{Scripts/iptables\_up.sh}:

\begin{lstlisting}[language=bash]
# Enable IP forwarding (required for FORWARD chain)
sudo sysctl -w net.ipv4.ip_forward=1

# Intercept forwarded traffic (IoT devices <-> Internet)
sudo iptables -I FORWARD -j NFQUEUE --queue-num 200 --queue-bypass

# Intercept traffic destined for the Pi itself
sudo iptables -I INPUT -j NFQUEUE --queue-num 100 --queue-bypass
\end{lstlisting}

\subsubsection{Teardown}

The cleanup script \texttt{Scripts/iptables\_down.sh} removes the NFQUEUE rules. The delete commands must match the insert commands exactly, including the \texttt{-{}-queue-bypass} flag; otherwise, \texttt{iptables} will not find the rule to delete. The script uses a \texttt{while} loop to handle duplicate rules that may exist if the setup script was run multiple times:

\begin{lstlisting}[language=bash]
# Delete rules (must include --queue-bypass to match)
while sudo iptables -D FORWARD -j NFQUEUE \
      --queue-num 200 --queue-bypass 2>/dev/null; do :; done
while sudo iptables -D INPUT -j NFQUEUE \
      --queue-num 100 --queue-bypass 2>/dev/null; do :; done
\end{lstlisting}

\subsubsection{NetfilterQueue Binding}

In Python, the \texttt{NetfilterQueue} library binds to the kernel queues and delivers raw packet payloads to a callback function. Two threads are created at startup, each binding to one queue number:

\begin{lstlisting}[language=python]
# Input chain thread
nfq = NetfilterQueue()
nfq.bind(100, callback)  # INPUT chain
nfq.run()

# Forward chain thread
nfq = NetfilterQueue()
nfq.bind(200, callback)  # FORWARD chain
nfq.run()
\end{lstlisting}

\subsection{Packet Parsing Module}
\label{sec:parsing}

The packet parser (\texttt{packet\_parser.py}) receives a raw NFQUEUE packet object and extracts the fields required by the detection engines. Since NFQUEUE delivers packets starting from the IP layer (no Ethernet header), the Scapy library is used to parse the IP, TCP, UDP, and ICMP layers.

\subsubsection{Extracted Fields}

Table~\ref{tab:parsed_fields} lists all fields extracted by the parser.

\begin{table}[H]
\centering
\caption{Fields extracted by the packet parser.}
\label{tab:parsed_fields}
\begin{tabular}{llp{7cm}}
\toprule
\textbf{Field} & \textbf{Type} & \textbf{Description} \\
\midrule
\texttt{src\_ip}     & \texttt{str}   & Source IP address \\
\texttt{dst\_ip}     & \texttt{str}   & Destination IP address \\
\texttt{src\_port}   & \texttt{int}   & Source port (0 if ICMP) \\
\texttt{dst\_port}   & \texttt{int}   & Destination port (0 if ICMP) \\
\texttt{port}        & \texttt{str}   & \texttt{"TCP"}, \texttt{"UDP"}, or \texttt{""} (for ICMP) \\
\texttt{tcp\_flags}  & \texttt{int}   & TCP flags bitmask (\texttt{SYN=0x02}, \texttt{ACK=0x10}, etc.) \\
\texttt{rawts}       & \texttt{float} & Raw Unix timestamp (seconds since epoch) \\
\texttt{timestamp}   & \texttt{str}   & Formatted timestamp (\texttt{YYYY-MM-DD HH:MM:SS}) \\
\texttt{packetID}    & \texttt{int}   & NFQUEUE packet identifier \\
\texttt{payloadLen}  & \texttt{int}   & Total IP payload length in bytes \\
\bottomrule
\end{tabular}
\end{table}

The parser returns a dictionary containing all these fields, which is consumed by both the behavioural and signature detection engines. Only TCP SYN packets (SYN flag set, ACK flag not set) are forwarded to the TCP analysis functions, as legitimate TCP traffic typically includes ACK flags in established connections.

\subsection{Behavioural Detection Engine}
\label{sec:behaviour}

The behavioural detection engine (\texttt{detectore\_engine.py}) implements anomaly-based detection for four attack types. It employs a dual-threshold strategy that combines a sliding time window with packet counting and an adaptive EWMA rate estimator that learns per-flow baselines from live traffic.

\subsubsection{Sliding Time Window}

For each flow key (e.g.\ a unique combination of destination IP and port), the engine maintains a \emph{deque} (double-ended queue) of packet timestamps. On each new packet, timestamps older than the configured window duration are removed from the front of the deque, and the new timestamp is appended to the back. The number of entries remaining in the deque represents the packet count within the current window.

This approach has $O(1)$ amortised time complexity per packet, since each timestamp is enqueued once and dequeued at most once. Memory usage is bounded by the maximum number of packets that can arrive within a single window period.

\subsubsection{EWMA Rate Estimator}

The EWMA estimator provides a smoothed, real-time measurement of the packet arrival rate in packets per second. Unlike the sliding window, which counts discrete packets, the EWMA captures the intensity of traffic and adapts smoothly to changes in rate. Each flow maintains its own independent estimator instance.

The fast rate update formula is:

\begin{equation}
    \Delta t = t_{\text{current}} - t_{\text{last}}
\end{equation}
\begin{equation}
    r_{\text{instant}} = \frac{1}{\Delta t}
\end{equation}
\begin{equation}
    r_{\text{ewma}} = \alpha \cdot r_{\text{instant}} + (1 - \alpha) \cdot r_{\text{ewma}}
\label{eq:ewma}
\end{equation}

The smoothing factor $\alpha = 0.3$ controls the trade-off between responsiveness and stability. The estimator reacts within a few packets to genuine bursts, whilst filtering out isolated spikes.

\subsubsection{Adaptive Threshold via Baseline Learning}

Rather than comparing the EWMA rate against a fixed, manually-chosen threshold, Loki IDS uses an adaptive threshold that learns from live traffic. Each flow maintains a second, much slower EWMA that tracks the ``normal'' baseline rate over time:

\begin{equation}
    r_{\text{baseline}} = \alpha_b \cdot r_{\text{instant}} + (1 - \alpha_b) \cdot r_{\text{baseline}}
\end{equation}
\begin{equation}
    T_{\text{adaptive}} = \max\!\left(r_{\text{baseline}} \times M,\; T_{\text{floor}}\right)
\end{equation}

The key idea is that each flow learns what its own normal traffic looks like, and the detection threshold adjusts accordingly. A flow that normally sees 10~pps will get a threshold of 50~pps (with the default multiplier $M=5$), whilst a flow that normally sees 2~pps will get a threshold of 10~pps. This eliminates the need to manually tune thresholds for each network environment.

\paragraph{Warm-Up Period.}
Before the baseline is considered reliable, the system requires a configurable number of warm-up packets (default: 50). During this warm-up phase, the $T_{\text{floor}}$ value is used as the threshold, providing a safe fallback so the system is still protected from attacks during initial startup.

\paragraph{Baseline Poisoning Protection.}
A critical concern with adaptive baselines is that an attacker could attempt to inflate the baseline by sending sustained high-rate traffic, thereby raising the threshold and hiding future attacks. To prevent this, the baseline EWMA \textbf{only updates when the current rate is below the current adaptive threshold}. If the traffic rate exceeds the threshold (i.e.\ a potential attack is in progress), the baseline learning is paused:

\begin{lstlisting}[language=python]
# Only learn from normal traffic, skip during attacks
if ewma_rate < current_adaptive_threshold:
    baseline_rate = alpha_baseline * instant_rate \
                  + (1 - alpha_baseline) * baseline_rate
\end{lstlisting}

This ensures attacks cannot poison the baseline.

\paragraph{Baseline Persistence Across Attacks.}
When an attack ends and the fast EWMA rate is reset, the baseline knowledge is preserved. The learned normal rate for a flow does not change simply because one attack ended. This means the system immediately resumes accurate detection after an attack without requiring another warm-up period.

\subsubsection{Dual-Threshold Detection Logic}

An alert is triggered only when \textbf{both} conditions are satisfied simultaneously:
\begin{enumerate}
    \item The sliding window packet count exceeds the count threshold (e.g.\ 200~packets in 2~seconds for TCP flood).
    \item The EWMA rate exceeds the adaptive threshold for that specific flow.
\end{enumerate}

This dual-check approach significantly reduces false positives. A brief burst of legitimate traffic might exceed one threshold but not both. Conversely, a sustained low-rate attack spread over a long period might not trigger the EWMA threshold even if it accumulates enough packets in the window. Only genuine flood attacks, characterised by both high volume and high rate, trigger detection.

\subsubsection{Detected Attack Types}

Table~\ref{tab:thresholds} summarises the detection parameters for each attack type.

\begin{table}[H]
\centering
\caption{Behavioural detection parameters with adaptive thresholds.}
\label{tab:thresholds}
\begin{tabular}{lllp{4.5cm}}
\toprule
\textbf{Attack Type} & \textbf{Window} & \textbf{Count Threshold} & \textbf{Adaptive EWMA} \\
\midrule
TCP SYN Port Scan & 5\,s & 20 unique ports   & N/A (count-only) \\
TCP SYN Flood     & 2\,s & 200 packets       & baseline $\times$ 5 (floor: 100\,pps) \\
UDP Flood         & 2\,s & 300 packets       & baseline $\times$ 5 (floor: 150\,pps) \\
ICMP Flood        & 2\,s & 100 packets       & baseline $\times$ 5 (floor: 50\,pps) \\
\bottomrule
\end{tabular}
\end{table}

Port scanning detection differs from flood detection in that it counts \emph{unique destination ports} rather than total packets. If a single source IP accesses more than 20 distinct ports on a destination IP within 5~seconds, a port scan alert is raised. EWMA is not applied to port scanning, as the distinguishing feature is port diversity rather than packet rate.

\subsection{Signature-Based Detection Engine}
\label{sec:signature}

The signature engine (\texttt{signature\_engine.py}) performs payload-based pattern matching against a database of known attack signatures. When a packet carries a Raw layer (application-layer data), its payload bytes are checked against each signature pattern.

\subsubsection{Signature Loading}

Signatures are loaded from the SQLite database at startup via the database integration module. Each signature consists of a name, a text pattern, an action, and a description. The text pattern is encoded to bytes at load time for efficient matching. Only enabled signatures are loaded:

\begin{lstlisting}[language=python]
# Signature format (from database)
{
    "name": "SQL Injection - Union Select",
    "pattern": "UNION SELECT",
    "pattern_bytes": b"UNION SELECT",  # pre-encoded
    "action": "alert",
    "description": "Detects UNION SELECT SQL injection"
}
\end{lstlisting}

\subsubsection{Matching Process}

For each packet with a Raw layer, the engine iterates through all loaded signatures and checks whether \texttt{pattern\_bytes in payload}. This is a simple substring search using Python's built-in \texttt{in} operator, which is implemented in C and runs efficiently. The first matching rule is returned.

\subsubsection{Example Signatures}

Table~\ref{tab:signatures} lists the default signature rules.

\begin{table}[H]
\centering
\caption{Example signature rules.}
\label{tab:signatures}
\begin{tabular}{llp{4cm}}
\toprule
\textbf{Name} & \textbf{Pattern} & \textbf{Category} \\
\midrule
SQL Injection -- Union Select  & \texttt{UNION SELECT}    & SQL Injection \\
SQL Injection -- OR 1=1        & \texttt{OR 1=1}          & SQL Injection \\
Path Traversal -- Etc Passwd   & \texttt{/etc/passwd}     & Path Traversal \\
XSS -- Script Tag              & \texttt{<script>}        & Cross-Site Scripting \\
XSS -- JavaScript Alert        & \texttt{javascript:alert}& Cross-Site Scripting \\
Command Injection -- Semicolon & \texttt{; /bin/}         & Command Injection \\
Git Repository Access          & \texttt{/.git/}          & Information Disclosure \\
Test Attack Signature          & \texttt{ATTACK\_TEST}    & Testing \\
\bottomrule
\end{tabular}
\end{table}

\subsection{Logging and Alert Lifecycle Management}
\label{sec:logging}

The logger module (\texttt{logger.py}) is responsible for recording detected alerts with smart deduplication. During a sustained attack, the IDS may detect the same attack hundreds or thousands of times per second. Without deduplication, this would flood the log files and database with identical entries.

\subsubsection{Alert Lifecycle}

Each unique alert progresses through a three-stage lifecycle:

\begin{description}[style=nextline]
    \item[STARTED] Logged immediately when a new attack is first detected. Contains the initial packet details.
    \item[ONGOING] Periodic updates logged during a sustained attack (at most 3~updates, every 5~seconds). Includes cumulative statistics: packet count, duration, and attack rate in packets per second.
    \item[ENDED] Logged when the attack ceases (no new packets for 10~seconds). Contains a comprehensive summary: total duration, total packets, average rate, first/last seen timestamps.
\end{description}

\subsubsection{Deduplication Logic}

Alerts are deduplicated by a composite key that groups related packets:
\begin{itemize}
    \item \textbf{Port scans:} key = \texttt{(type, message, src\_ip, dst\_ip, "scan")}.
    \item \textbf{Floods:} key = \texttt{(type, message, src\_ip, dst\_ip, dst\_port)}.
\end{itemize}
All subsequent packets matching an existing key update the internal state without generating new log entries, except for the periodic ONGOING updates.

\subsubsection{Alert Types and Subtypes}

Table~\ref{tab:alert_types} shows the alert classification taxonomy.

\begin{table}[H]
\centering
\caption{Alert classification taxonomy.}
\label{tab:alert_types}
\begin{tabular}{llp{5.5cm}}
\toprule
\textbf{Type} & \textbf{Subtype} & \textbf{Description} \\
\midrule
BEHAVIOR  & PORT\_SCAN  & TCP port scanning detection \\
BEHAVIOR  & TCP\_FLOOD  & TCP SYN flood (DoS/DDoS) \\
BEHAVIOR  & UDP\_FLOOD  & UDP flood (DoS/DDoS) \\
BEHAVIOR  & ICMP\_FLOOD & ICMP flood (DoS/DDoS) \\
SIGNATURE & (none)      & Payload pattern match \\
SYSTEM    & (none)      & System events (startup, shutdown, errors) \\
\bottomrule
\end{tabular}
\end{table}

\subsubsection{Output Formats}

All alerts are written to a JSONL (JSON Lines) file at \texttt{logs/loki\_alerts.jsonl}, with one JSON object per line. Each record contains a timestamp, alert status, type, subtype, source/destination IP and port, a human-readable message, and a details object with additional context. When database integration is enabled, alerts are concurrently written to the SQLite database.

Example JSONL record (STARTED):

\begin{lstlisting}[language=json]
{
  "timestamp": "2026-01-30T22:15:00.123456",
  "status": "STARTED",
  "type": "BEHAVIOR",
  "subtype": "TCP_FLOOD",
  "src_ip": "10.0.0.50",
  "dst_ip": "10.0.0.1",
  "src_port": 12345,
  "dst_port": 443,
  "message": "TCP Flood (DoS/DDoS) Detected on FORWARD chain",
  "details": {"dst_ip": "10.0.0.1", "chain": "FORWARD"}
}
\end{lstlisting}

\subsection{Database Integration}
\label{sec:db_integration}

The database integration module (\texttt{db\_integration.py}) provides a bridge between the synchronous IDS packet-processing loop and the asynchronous SQLAlchemy database layer used by the web interface. It exposes two primary operations: inserting alerts and retrieving signature rules.

Since the IDS runs in a synchronous context (NFQUEUE callbacks are not async), the module spawns background threads with their own \texttt{asyncio} event loops to perform database operations without blocking packet processing. Alert insertions are fire-and-forget (non-blocking), whilst signature retrieval is blocking (the IDS waits for signatures to be loaded before starting detection).

The database uses SQLite with the \texttt{aiosqlite} driver, stored at \texttt{Web-Interface/loki\_ids.db}. The schema includes tables for alerts, signatures, and a statistics cache.

\subsection{Configuration and Thresholds}
\label{sec:config}

All detection thresholds are configurable via the \texttt{PortScanningDetector} constructor in \texttt{detectore\_engine.py}. The following tables summarise all tuneable parameters.

\subsubsection{Sliding Window Parameters}

\begin{table}[H]
\centering
\caption{Sliding window detection thresholds.}
\label{tab:window_params}
\begin{tabular}{llll}
\toprule
\textbf{Parameter} & \textbf{Default} & \textbf{Unit} & \textbf{Description} \\
\midrule
\texttt{port\_scanning\_window}    & 5   & seconds      & Window duration for port scan detection \\
\texttt{port\_scanning\_threshold} & 20  & unique ports & Distinct ports to trigger alert \\
\texttt{tcp\_flood\_window}        & 2   & seconds      & Window duration for TCP flood \\
\texttt{tcp\_flood\_threshold}     & 200 & packets      & Packet count to trigger TCP flood \\
\texttt{udp\_flood\_window}        & 2   & seconds      & Window duration for UDP flood \\
\texttt{udp\_flood\_threshold}     & 300 & packets      & Packet count to trigger UDP flood \\
\texttt{icmp\_flood\_window}       & 2   & seconds      & Window duration for ICMP flood \\
\texttt{icmp\_flood\_threshold}    & 100 & packets      & Packet count to trigger ICMP flood \\
\bottomrule
\end{tabular}
\end{table}

\subsubsection{Adaptive EWMA Parameters}

\begin{table}[H]
\centering
\caption{Adaptive EWMA threshold parameters.}
\label{tab:ewma_params}
\begin{tabular}{llp{6.5cm}}
\toprule
\textbf{Parameter} & \textbf{Default} & \textbf{Description} \\
\midrule
\texttt{alpha}              & 0.3      & Fast EWMA smoothing factor. Tracks current rate. Higher = more reactive to bursts. \\
\texttt{alpha\_baseline}    & 0.05     & Slow EWMA smoothing factor. Learns baseline ``normal'' rate. Much smaller so it moves slowly. \\
\texttt{multiplier}         & 5.0      & How many times above baseline constitutes an attack. $T = \text{baseline} \times \text{multiplier}$. \\
\texttt{minimum\_floor} (TCP)  & 100.0\,pps  & Minimum threshold floor for TCP flood. Used during warm-up. \\
\texttt{minimum\_floor} (UDP)  & 150.0\,pps  & Minimum threshold floor for UDP flood. \\
\texttt{minimum\_floor} (ICMP) & 50.0\,pps   & Minimum threshold floor for ICMP flood. \\
\texttt{warmup\_packets}    & 50       & Packets required before baseline is trusted. During warm-up, $T_{\text{floor}}$ is used. \\
\bottomrule
\end{tabular}
\end{table}

\subsubsection{Logger Parameters}

\begin{table}[H]
\centering
\caption{Logger configuration parameters.}
\label{tab:logger_params}
\begin{tabular}{llp{6.5cm}}
\toprule
\textbf{Parameter} & \textbf{Default} & \textbf{Description} \\
\midrule
\texttt{alert\_cooldown}   & 10\,s & Time of inactivity before an attack is considered ended \\
\texttt{update\_interval}  & 5\,s  & Minimum interval between ONGOING log entries \\
\texttt{max\_updates}      & 3     & Maximum number of ONGOING entries per attack \\
\bottomrule
\end{tabular}
\end{table}

\subsection{API Server}
\label{sec:api}

The API server is built with FastAPI and serves both the REST API and the static web dashboard.

\begin{table}[H]
\centering
\caption{Key API endpoints.}
\label{tab:endpoints}
\begin{tabular}{lll}
\toprule
\textbf{Method} & \textbf{Endpoint} & \textbf{Purpose} \\
\midrule
GET    & \texttt{/api/alerts}                     & List alerts (filtered, paginated) \\
POST   & \texttt{/api/alerts}                     & Create alert (from IDS core) \\
GET    & \texttt{/api/signatures}                 & List signatures \\
POST   & \texttt{/api/signatures}                 & Create signature \\
PUT    & \texttt{/api/signatures/\{id\}}          & Update signature \\
DELETE & \texttt{/api/signatures/\{id\}}          & Delete signature \\
POST   & \texttt{/api/signatures/reload}          & Import from YAML \\
GET    & \texttt{/api/stats}                      & Aggregated statistics \\
GET    & \texttt{/api/system/health}              & Health check \\
GET    & \texttt{/api/system/status}              & IDS running status \\
GET    & \texttt{/api/iot/devices}                & List IoT devices \\
POST   & \texttt{/api/iot/devices/\{id\}/alarm}   & Control alarm (ESP32-1) \\
POST   & \texttt{/api/iot/devices/\{id\}/bulb}    & Control bulb (ESP32-2) \\
WS     & \texttt{/ws/alerts}                      & Real-time alert streaming \\
\bottomrule
\end{tabular}
\end{table}

The database uses SQLite with WAL (Write-Ahead Logging) mode enabled for concurrent read/write access. The async engine (via \texttt{aiosqlite}) ensures that database operations do not block the event loop.

\subsection{Project File Structure}
\label{sec:structure}

\begin{lstlisting}[frame=none,backgroundcolor=\color{white},basicstyle=\ttfamily\footnotesize]
Loki-IDS/
|-- Core/
|   +-- loki/
|       |-- nfqueue_app.py          # Main entry point
|       |-- detectore_engine.py     # Behavioural detection
|       |                            (sliding window +
|       |                             adaptive EWMA)
|       |-- signature_engine.py     # Signature pattern matching
|       |-- packet_parser.py        # Scapy packet field extraction
|       |-- logger.py               # Alert logging + deduplication
|       |-- db_integration.py       # Async DB bridge
|       +-- api/
|           +-- main.py             # FastAPI application
|           +-- models/             # SQLAlchemy models, schemas, CRUD
|           +-- routes/             # API route handlers
|           +-- iot/                # MQTT client (paho-mqtt 2.x)
|
|-- Web-Interface/static/           # Frontend (HTML, JS, CSS)
|
|-- Scripts/
|   |-- install_deps.sh             # Dependency installer
|   |-- iptables_up.sh              # Setup NFQUEUE rules
|   +-- iptables_down.sh            # Remove NFQUEUE rules
|
|-- attack-scripts/
|   |-- nmap-portscan.sh            # Port scan test
|   |-- dos-tcp.sh                  # TCP flood test
|   |-- dos-udp.sh                  # UDP flood test
|   +-- icmp-attack.sh              # ICMP flood test
|
|-- Configs/
|   +-- test_yaml_file.yaml         # Legacy YAML signatures
|
|-- logs/
|   +-- loki_alerts.jsonl           # Alert log output
|
+-- run_loki.sh                     # One-command launcher
+-- start_loki_system.sh            # Full system startup
+-- setup_venv.sh                   # Virtual environment setup
+-- requirements.txt                # Python dependencies
\end{lstlisting}

\subsection{Testing}
\label{sec:testing}

The IDS was tested against all supported attack categories to verify correct detection, alert lifecycle management, and dashboard integration. All attacks were launched from a dedicated attacker Raspberry~Pi connected to the network through the MikroTik router, simulating an external adversary who is not a member of the IoT Wi-Fi network.

\subsubsection{Attack Scripts}
\label{sec:attack_scripts}

To facilitate reproducible testing, a set of attack scripts is provided in the \texttt{attack-scripts/} directory of the project repository. These scripts are deployed on the attacker Raspberry~Pi's home directory and accept the target IP address as a command-line argument. Table~\ref{tab:attack_scripts} summarises the available scripts.

\begin{table}[H]
\centering
\caption{Attack scripts for IDS testing.}
\label{tab:attack_scripts}
\begin{tabular}{llp{5.5cm}}
\toprule
\textbf{Script} & \textbf{Tool} & \textbf{Purpose} \\
\midrule
\texttt{nmap-portscan.sh} & Nmap   & TCP service/version scan of the top 1000 ports \\
\texttt{dos-tcp.sh}       & Nping  & TCP flood: sends 1000 TCP packets to port 10001 \\
\texttt{dos-udp.sh}       & Nping  & UDP flood: sends 1000 UDP packets to port 10001 \\
\texttt{icmp-attack.sh}   & Nping  & ICMP flood: sends 1000 ICMP echo requests \\
\bottomrule
\end{tabular}
\end{table}

The contents of each script are shown below.

\paragraph{Port Scan (\texttt{attack-scripts/nmap-portscan.sh}).}
\begin{lstlisting}[language=bash]
sudo nmap -sV --top-ports 1000 $1
\end{lstlisting}

This script performs a service/version scan (\texttt{-sV}) of the 1000 most commonly used TCP ports on the target. Nmap sends SYN probes to each port and, for open ports, attempts to identify the running service. The IDS detects this as a port scan because the attacker contacts more than 20 distinct destination ports within the 5-second sliding window.

\paragraph{TCP Flood (\texttt{attack-scripts/dos-tcp.sh}).}
\begin{lstlisting}[language=bash]
#!/bin/bash
sudo nping -tcp -p 10001 -c 1000 $1
\end{lstlisting}

Sends 1000 TCP packets to port 10001 at the maximum rate permitted by Nping. The high packet rate triggers both the sliding window count threshold (200 packets in 2~seconds) and the adaptive EWMA rate threshold.

\paragraph{UDP Flood (\texttt{attack-scripts/dos-udp.sh}).}
\begin{lstlisting}[language=bash]
sudo nping -udp -p 10001 -c 1000 $1
\end{lstlisting}

Sends 1000 UDP datagrams to port 10001. The detection logic mirrors the TCP flood case but uses the UDP-specific thresholds (300 packets in 2~seconds, EWMA floor of 150\,pps).

\paragraph{ICMP Flood (\texttt{attack-scripts/icmp-attack.sh}).}
\begin{lstlisting}[language=bash]
#!/bin/bash
sudo nping -icmp -c 1000 $1
\end{lstlisting}

Sends 1000 ICMP echo request packets. The ICMP flood detector triggers at 100 packets within the 2-second window combined with the EWMA rate exceeding 50\,pps.

\paragraph{Usage.} Each script is executed from the attacker Raspberry~Pi with the target IP address as the sole argument:

\begin{lstlisting}[language=bash]
# From the attacker Raspberry Pi
bash dos-tcp.sh 10.0.0.1
\end{lstlisting}

\subsubsection{Test Results}

\paragraph{1. Port Scan Detection.}

\begin{lstlisting}[language=bash]
# From attacker RPi
bash nmap-portscan.sh 10.0.0.1
\end{lstlisting}

\textbf{Result:} Loki detected the scan within 2~seconds. The dashboard showed a BEHAVIOR alert with subtype PORT\_SCAN, correctly identifying the source IP and listing the scanned ports. The alert lifecycle progressed from STARTED through ONGOING to ENDED.

\paragraph{2. TCP Flood.}

\begin{lstlisting}[language=bash]
# From attacker RPi
bash dos-tcp.sh 10.0.0.1
\end{lstlisting}

\textbf{Result:} Alert raised after the EWMA rate exceeded 100\,pps and the sliding window count exceeded 200. ONGOING updates showed the attack rate climbing to several thousand packets per second. The ENDED alert reported the total packet count and duration.

\paragraph{3. UDP Flood.}

\begin{lstlisting}[language=bash]
# From attacker RPi
bash dos-udp.sh 10.0.0.1
\end{lstlisting}

\textbf{Result:} UDP\_FLOOD alert raised with correct thresholds. The dual-check mechanism prevented false positives from normal DNS traffic.

\paragraph{4. ICMP Flood.}

\begin{lstlisting}[language=bash]
# From attacker RPi
bash icmp-attack.sh 10.0.0.1
\end{lstlisting}

\textbf{Result:} ICMP\_FLOOD alert raised correctly. The EWMA threshold of 50\,pps was reached quickly during the flood.

\paragraph{5. Signature Detection.}

\begin{lstlisting}[language=bash]
# From attacker RPi
echo "ATTACK_TEST" | nc 10.0.0.1 80
\end{lstlisting}

\textbf{Result:} SIGNATURE alert with the matching pattern ``ATTACK\_TEST'' displayed in the dashboard. The source IP and matched rule name were correctly recorded.

\paragraph{6. Normal Traffic (False Positive Test).}

Regular web browsing, MQTT communication between ESP32 devices, and SSH sessions were tested to verify they did not trigger alerts. No false positives were observed during normal operation.

\subsubsection{Supervisor Testing Procedure}
\label{sec:supervisor_testing}

The following step-by-step procedure allows the supervisor to independently verify the IDS functionality on the prepared test environment. All equipment is pre-configured; no additional setup is required.

\paragraph{Step 1: Verify the Environment.}
\begin{enumerate}
    \item Power on the Raspberry~Pi gateway. It will automatically start the access point and begin broadcasting the \texttt{Access-Point} Wi-Fi network.
    \item Confirm that the MikroTik router is powered on and connected to the Pi's Ethernet port (\texttt{eth0}).
    \item Connect the attacker Raspberry~Pi (or a laptop) to the MikroTik router via Ethernet or the router's Wi-Fi.
\end{enumerate}

\paragraph{Step 2: Start the Loki IDS System.}
On the Raspberry~Pi gateway, start the full system:
\begin{lstlisting}[language=bash]
cd ~/Loki-IDS
sudo bash start_loki_system.sh
\end{lstlisting}
This script starts the API server on port 8080, configures the iptables NFQUEUE rules, and launches the IDS engine. Wait until the terminal output confirms that both NFQUEUE threads (INPUT queue~100 and FORWARD queue~200) are running.

\paragraph{Step 3: Access the Dashboard.}
Open a web browser on a device connected to the access point and navigate to:
\begin{lstlisting}[frame=none,backgroundcolor=\color{white}]
http://10.0.0.1:8080
\end{lstlisting}
Verify that the dashboard loads and the \textbf{Dashboard} tab shows the system status as ``Running''.

\paragraph{Step 4: Verify IoT Device Connectivity.}
\begin{enumerate}
    \item Power on both ESP32 modules. They will automatically connect to the access point and the MQTT broker.
    \item In the dashboard, navigate to the \textbf{IoT Control} tab and confirm that both devices appear as ``Online'' (indicated by periodic MQTT heartbeat messages).
    \item Test device control: toggle the alarm on ESP32-1 and adjust the bulb brightness on ESP32-2 from the dashboard to confirm end-to-end MQTT communication.
\end{enumerate}

\paragraph{Step 5: Run Attack Scripts from the Attacker Machine.}
From the attacker Raspberry~Pi (connected via the MikroTik router), execute each attack script in sequence, allowing approximately 15~seconds between scripts for the previous alert to reach its ENDED state:

\begin{lstlisting}[language=bash]
# 1. Port Scan
bash ~/nmap-portscan.sh 10.0.0.1

# 2. TCP Flood
bash ~/dos-tcp.sh 10.0.0.1

# 3. UDP Flood
bash ~/dos-udp.sh 10.0.0.1

# 4. ICMP Flood
bash ~/icmp-attack.sh 10.0.0.1
\end{lstlisting}

After each script completes, switch to the dashboard's \textbf{Alerts} tab and verify:
\begin{itemize}
    \item A new alert has appeared with the correct type (BEHAVIOR) and subtype (PORT\_SCAN, TCP\_FLOOD, UDP\_FLOOD, or ICMP\_FLOOD).
    \item The source IP matches the attacker machine's IP address.
    \item The alert lifecycle transitions from STARTED to ONGOING and finally to ENDED after approximately 10~seconds of inactivity.
    \item A real-time toast notification appeared when the alert was first raised.
\end{itemize}

\paragraph{Step 6: Test Signature Detection.}
\begin{lstlisting}[language=bash]
# From the attacker machine
echo "ATTACK_TEST" | nc 10.0.0.1 80
\end{lstlisting}
Verify that a SIGNATURE alert appears in the dashboard with the matched rule name ``Test Attack Signature''.

\paragraph{Step 7: Inspect Alert Logs.}
On the Raspberry~Pi gateway, inspect the alert log file to confirm all alerts were persisted:
\begin{lstlisting}[language=bash]
cat ~/Loki-IDS/logs/loki_alerts.jsonl | python3 -m json.tool
\end{lstlisting}
Each attack should have produced three log entries: STARTED, ONGOING, and ENDED.

\paragraph{Step 8: Verify Normal Traffic (No False Positives).}
From a device connected to the access point, perform normal operations---browse the dashboard, let ESP32 devices send heartbeats, and initiate SSH sessions to the Pi. Confirm that no BEHAVIOR or SIGNATURE alerts are raised during this period.

\paragraph{Step 9: Shutdown.}
\begin{lstlisting}[language=bash]
# Press Ctrl+C in the IDS terminal, then:
sudo bash ~/Loki-IDS/Scripts/iptables_down.sh
\end{lstlisting}
The IDS performs a graceful shutdown: all active alerts are finalised with ENDED status and iptables rules are removed.

\subsection{Deployment and Execution}
\label{sec:deployment}

\subsubsection*{Prerequisites}
\begin{itemize}
    \item Raspberry~Pi (or any Linux system) running Debian/Ubuntu with root access.
    \item Python 3.12 or later (Python 3.13 recommended).
    \item System packages: \texttt{build-essential}, \texttt{python3-dev}, \texttt{python3-venv}, \texttt{libnetfilter-queue-dev}, \texttt{libnfnetlink-dev}, \texttt{libpcap-dev}.
    \item Python packages: \texttt{NetfilterQueue}, \texttt{scapy}, \texttt{pyyaml}.
\end{itemize}

\subsubsection*{Quick Start (One Command)}

The simplest way to run the IDS is to use the launcher script from the project root:

\begin{lstlisting}[language=bash]
cd /path/to/Loki-IDS
sudo bash run_loki.sh
\end{lstlisting}

This script automatically sets up iptables rules, activates the virtual environment, starts the IDS, and cleans up iptables rules when \texttt{Ctrl+C} is pressed.

\subsection{Reproducible Setup Steps}
\label{sec:ids_steps}

This section provides a complete, reproducible procedure for deploying Loki IDS on a Raspberry~Pi from a clean installation.

\subsubsection*{Step 1: Clone the Repository}

\begin{lstlisting}[language=bash]
cd ~
git clone <repository-url> Loki-IDS
cd Loki-IDS
\end{lstlisting}

\subsubsection*{Step 2: Install System Dependencies}

\begin{lstlisting}[language=bash]
sudo apt update
sudo apt install -y build-essential python3-dev python3-venv \
    libnetfilter-queue-dev libnfnetlink-dev libpcap-dev
\end{lstlisting}

\subsubsection*{Step 3: Create the Python Virtual Environment}

\begin{lstlisting}[language=bash]
python3 -m venv loki_env
source loki_env/bin/activate
pip install --upgrade pip
pip install NetfilterQueue scapy pyyaml
\end{lstlisting}

Alternatively, run \texttt{bash Scripts/install\_deps.sh} which performs steps~2 and~3 automatically.

\subsubsection*{Step 4: Initialise the Database (Optional)}

If using the web interface for signature management and alert storage:

\begin{lstlisting}[language=bash]
cd Web-Interface
venv/bin/python3 init_database.py
cd ..
\end{lstlisting}

\subsubsection*{Step 5: Configure iptables}

\begin{lstlisting}[language=bash]
sudo bash Scripts/iptables_up.sh
\end{lstlisting}

Verify that the NFQUEUE rules are active:

\begin{lstlisting}[language=bash]
sudo iptables -L --line-numbers | grep NFQUEUE
# Expected output:
# 1  NFQUEUE  ...  NFQUEUE num 100 bypass
# 1  NFQUEUE  ...  NFQUEUE num 200 bypass
\end{lstlisting}

\subsubsection*{Step 6: Start the IDS}

\begin{lstlisting}[language=bash]
source loki_env/bin/activate
cd Core/loki
sudo python3 nfqueue_app.py
\end{lstlisting}

The IDS will output startup messages confirming database integration, signature loading, and thread initialisation. Packet-level logs will appear in real time. During the initial warm-up period (first $\sim$50~packets per flow), the adaptive EWMA thresholds will use their minimum floor values. After warm-up completes, the thresholds will automatically adjust to the learned baseline for each flow.

\subsubsection*{Step 7: Verify Detection}

From the attacker Raspberry~Pi (connected via the MikroTik router), run the provided attack scripts to verify detection is working:

\begin{lstlisting}[language=bash]
# Port scan (from attacker RPi)
bash ~/nmap-portscan.sh 10.0.0.1

# TCP flood
bash ~/dos-tcp.sh 10.0.0.1

# UDP flood
bash ~/dos-udp.sh 10.0.0.1

# ICMP flood
bash ~/icmp-attack.sh 10.0.0.1

# Signature test (send payload containing pattern)
echo "ATTACK_TEST" | nc 10.0.0.1 80
\end{lstlisting}

\subsubsection*{Step 8: Stop and Clean Up}

\begin{lstlisting}[language=bash]
# Press Ctrl+C in the IDS terminal, then:
sudo bash Scripts/iptables_down.sh
deactivate
\end{lstlisting}

The system performs a graceful shutdown: all active alerts are finalised with ENDED status, session statistics are logged, and (if using \texttt{run\_loki.sh}) iptables rules are automatically removed.

\subsubsection*{Step 9: Review Logs}

\begin{lstlisting}[language=bash]
# View alerts
cat logs/loki_alerts.jsonl | python3 -m json.tool

# Count alerts by type
grep -c '"type": "BEHAVIOR"' logs/loki_alerts.jsonl
grep -c '"type": "SIGNATURE"' logs/loki_alerts.jsonl
\end{lstlisting}


% ══════════════════════════════════════════════════════════
% REFERENCES
% ══════════════════════════════════════════════════════════
\newpage
\begin{thebibliography}{14}

\bibitem{atzori2010iot}
L.~Atzori, A.~Iera, and G.~Morabito,
``The Internet of Things: A survey,''
\emph{Computer Networks}, vol.~54, no.~15, pp.~2787--2805, 2010.

\bibitem{alfuqaha2015iot}
A.~Al-Fuqaha, M.~Guizani, M.~Mohammadi, M.~Aledhari, and M.~Ayyash,
``Internet of Things: A survey on enabling technologies, protocols, and applications,''
\emph{IEEE Communications Surveys \& Tutorials}, vol.~17, no.~4, pp.~2347--2376, 2015.

\bibitem{khan2012future}
R.~Khan, S.~U.~Khan, R.~Zaheer, and S.~Khan,
``Future Internet: The Internet of Things architecture, possible applications and key challenges,''
in \emph{Proc.\ 10th International Conference on Frontiers of Information Technology}, 2012, pp.~257--260.

\bibitem{wu2010architecture}
M.~Wu, T.-J.~Lu, F.-Y.~Ling, J.~Sun, and H.-Y.~Du,
``Research on the architecture of Internet of Things,''
in \emph{Proc.\ 3rd International Conference on Advanced Computer Theory and Engineering}, 2010, vol.~5, pp.~484--487.

\bibitem{banks2014mqtt}
A.~Banks and R.~Gupta,
``MQTT Version 3.1.1,''
OASIS Standard, 2014.

\bibitem{owasp2018iot}
OWASP,
``OWASP Internet of Things Top 10,''
2018.
[Online]. Available: \url{https://owasp.org/www-project-internet-of-things/}

\bibitem{bertino2017botnets}
E.~Bertino and N.~Islam,
``Botnets and Internet of Things Security,''
\emph{Computer}, vol.~50, no.~2, pp.~76--79, 2017.

\bibitem{ferrag2020deep}
M.~A.~Ferrag, L.~Maglaras, S.~Moschoyiannis, and H.~Janicke,
``Deep learning for cyber security intrusion detection: Approaches, datasets, and comparative study,''
\emph{Journal of Information Security and Applications}, vol.~50, p.~102419, 2020.

\bibitem{neshenko2019demystifying}
N.~Neshenko, E.~Bou-Harb, J.~Crichigno, G.~Kaddoum, and N.~Ghani,
``Demystifying IoT security: An exhaustive survey on IoT vulnerabilities and a first empirical look on Internet-scale IoT exploitations,''
\emph{IEEE Communications Surveys \& Tutorials}, vol.~21, no.~3, pp.~2702--2733, 2019.

\bibitem{antonakakis2017mirai}
M.~Antonakakis \emph{et~al.},
``Understanding the Mirai Botnet,''
in \emph{Proc.\ 26th USENIX Security Symposium}, 2017, pp.~1093--1110.

\bibitem{khraisat2019survey}
A.~Khraisat, I.~Gondal, P.~Vamplew, and J.~Kamruzzaman,
``Survey of intrusion detection systems: techniques, datasets and challenges,''
\emph{Cybersecurity}, vol.~2, no.~1, pp.~1--22, 2019.

\bibitem{sommer2010outside}
R.~Sommer and V.~Paxson,
``Outside the closed world: On using machine learning for network intrusion detection,''
in \emph{Proc.\ IEEE Symposium on Security and Privacy}, 2010, pp.~305--316.

\bibitem{alaiz2019multiclass}
H.~Alaiz-Moreton \emph{et~al.},
``Multiclass classification procedure for detecting attacks on MQTT-IoT protocol,''
\emph{Complexity}, vol.~2019, pp.~1--11, 2019.

\bibitem{page2021prisma}
M.~J.~Page \emph{et~al.},
``The PRISMA 2020 statement: An updated guideline for reporting systematic reviews,''
\emph{BMJ}, vol.~372, p.~n71, 2021.

\end{thebibliography}

\end{document}
